% Generated draft: requires mathematical and source review.
\documentclass[11pt]{article}
\usepackage[margin=25mm]{geometry}
\usepackage{fontspec}
\setmainfont{Latin Modern Roman}
\usepackage{amsmath,amssymb,mathtools,unicode-math}
\setmathfont{Latin Modern Math}
\usepackage{enumitem,needspace,xurl,hyperref}
\setlength{\parindent}{0pt}
\setlength{\parskip}{5pt}
\setlength{\emergencystretch}{3em}
\newcommand{\N}{\mathbb{N}}
\newcommand{\Z}{\mathbb{Z}}
\newcommand{\Q}{\mathbb{Q}}
\newcommand{\R}{\mathbb{R}}
\newcommand{\C}{\mathbb{C}}
\newcommand{\F}{\mathbb{F}}
\newcommand{\A}{\mathbb{A}}
\newcommand{\PP}{\mathbb{P}}
\newcommand{\E}{\mathbb{E}}
\newcommand{\eps}{\varepsilon}
\newcommand{\dd}{\,\mathrm{d}}
\begin{document}
\section*{UnsolvedMath ID: 1900013}
\textbf{Source problem code:} catalogue entry
\subsection{Definitions and mathematical statement}
A cubic extension of $\Q$ is a field $K\supset\Q$ with $[K:\Q]=3$. Two such fields are the same up to $\Q$-isomorphism if there is a field isomorphism between them fixing $\Q$.

\textbf{Problem attributed to V. Arnold.} Classify the continued fractions corresponding to the same cubic extension of $\Q$. The title places this question in the combinatorial structure of sails in multidimensional continued-fraction geometry. The record does not define the continued-fraction construction, the correspondence with a cubic field, or the equivalence relation under which a classification is requested; ordinary one-dimensional continued fractions are not silently substituted for the sail objects.

One standard source of cubic-field data in this subject is an integer matrix $A\in M_3(\Z)$ with irreducible characteristic polynomial of degree $3$: then
\[
 \Q[A]\cong\Q[t]/(\det(tI-A))
\]
is a cubic field. In a totally real eigenstructure, associated cones and sails can supply geometric continued-fraction data. This explains a possible connection with the problem, but the supplied excerpt does not identify this as its precise construction, restrict to totally real fields, or impose a condition such as $\det A=\pm1$.

Once the missing construction is fixed, let $\mathcal C$ denote its continued-fraction objects and let $\kappa(C)$ denote the associated cubic field. For every cubic field $K$ in the intended scope, the target is the classification of
\[
 \{C\in\mathcal C:\kappa(C)\cong_{\Q}K\}
\]
under the intended equivalence, which could distinguish combinatorial face-incidence type from integral-linear equivalence or other arithmetic data. The source must specify these choices before this becomes a unique classification problem.

The supplied research summary labels the problem OPEN-TRIAGE with moderate confidence, reports that no complete classification over a fixed cubic field was found, and says that the periodic--Markov structure is partially understood. It gives neither a definition of that structure nor citations for the partial results. These are recorded status claims rather than verified conclusions.
\subsection{Short English statement}
Arnold asks for a classification of the multidimensional continued fractions associated with one cubic number field. The record places the problem in sail geometry but omits the field correspondence and the intended notion of equivalence, so a unique precise classification target cannot yet be recovered.
\subsection{Sources}
\begin{itemize}
\item Local UnsolvedMath registry, permanent ID 1900013, source code catalogue entry, ``Geometry of Continued Fractions --- Combinatorial structure of sails''; attribution in the statement: V. Arnold. Supplied status open, classification OPEN-TRIAGE, updated 2026-08-06.
\item Karpenkov, \emph{Open problems in geometry of continued fractions} (2017), Problem 13, \url{https://arxiv.org/abs/1712.01450}. The record reports source-TeX extraction and access on 2026-07-29. It says the source presents an open item and current status requires release review. Rights: NEEDS\_REVIEW; public arXiv source TeX inspected according to the record.
\item The supplied research summary describes moderate-confidence openness and unspecified partial understanding of a periodic--Markov structure; no separate reference for these claims is supplied.
\end{itemize}
\end{document}
